\section{Related work and parameter comparison}
\label{sec:prior-art}

\subsection{Quantum algorithms for other Ramsey tasks}

Several quantum approaches to Ramsey problems address a different input and
output task from the one studied here.  Work based on adiabatic computation,
quantum annealing, or quantum counting searches for colourings that establish
or test Ramsey-number bounds
\cite{GaitanClark2012,BianEtAl2013,Wang2016,QuLiWangBaoCao2013,
RanjbarMacreadyClarkGaitan2016,PionMniszewski2025}.  In those formulations,
the search space consists of labelled graph colourings.  Our input is one
fixed graph given by an adjacency oracle, and the goal is to find a
homogeneous set guaranteed to exist in that graph.

Generic quantum algorithms for clique, independent set, and fixed-subgraph
containment operate in a closer oracle model
\cite{Doern2005,ChildsEisenberg2005,LeeMagniezSantha2012,Zhu2012}.  Quantum
backtracking supplies another general search primitive \cite{Montanaro2018}.
Our algorithms additionally exploit the nested majority structure of the
constructive Ramsey proof: each candidate set is represented by the
constraints accumulated along the recursion and sampled as an implicit set.
We are not aware of an earlier implementation of the constructive
Erd\H{o}s--Szekeres recursion by this mechanism in the adjacency-query model.

\subsection{Constructive Ramsey search}

Ramsey search also appears in proof complexity and TFNP.  It is connected to
weak and structured pigeonhole principles
\cite{Krajicek2001,Krajicek2005}, its white-box hardness follows from
collision-resistant hashing \cite{KomargodskiNaorYogev2019}, and it belongs
to the polynomial long-choice class PLC
\cite{PasarkarPapadimitriouYannakakis2023}.  The closest matching query
relation is studied in \cite{JainLiRobereXun2024}, which shows that it is
not black-box reducible to PIGEON and, in the paragraph ``Query complexity
of RAMSEY'' of its Section~III, raises the question of its query
complexity in the deterministic, randomized, and quantum models.  On the
classical side, the online Ramsey game (see \cite{ConlonFoxGrinshpunHe2019}
and the references therein), in which Builder queries edges and Painter
answers adversarially, is exactly the deterministic version of this
question with an unbounded vertex set; the best known lower bound for it is
the random-Painter bound of \cite{ConlonFoxGrinshpunHe2019} used in
\cref{prop:classical-lower}.

\subsection{Comparison with a reported quantum lower bound}
\label{sec:jlrx}

\paragraph{Model and parameter alignment.}
Definition~II.6 of \cite{JainLiRobereXun2024} uses $N=2^n$ vertices, and the
convention following Definition~II.7 sets the default target to $K=n/2$.
The unnumbered ``Query complexity of RAMSEY'' paragraph in Section~III
(p.~415) reports an $N^{1-o(1)}$ quantum lower bound, based on its
Lemma~III.1 and the multicollision lower bound of \cite{LiuZhandry2019}.

Away from the diagonal, the graph produced by Lemma~III.1 from a function
$h$ and a fixed graph $A_0$ is
\[
 A_h(u,v)=\mathbf 1[h(u)=h(v)]\lor A_0(h(u),h(v)),
 \qquad A_h(u,u)=0.
\]
It is a symmetric simple graph.  A coherent query uses $O(1)$ queries to the
value oracle for $h$, including uncomputation of the two hash values.  Thus
the reduction lies within our valid-graph promise.  Moreover,
\cref{cor:unpromised} gives the same upper bound for the locally verifiable
arbitrary-circuit relation, and \cref{cor:succinct} returns
$\lfloor n/2\rfloor+1$ vertices, from which one may discard a vertex to meet
the default target.  The graph promise, encoding, and output size therefore
align in the two formulations.

\paragraph{Exponent obtained by direct substitution.}
Let $H$ denote the multicollision range size and $G$ the number of vertices in
the resulting Ramsey instance.  For fixed multiplicity $t$, the sufficient
condition in Theorem~I.3 of \cite{JainLiRobereXun2024} is
$G\ge H^{4t}/4^t$.  At the boundary $G=\Theta_t(H^{4t})$, the exponent in
\cite{LiuZhandry2019} is
\[
 \alpha_t=\frac{2^{t-1}-1}{2^t-1}.
\]
The exponent transferred to $G$ is therefore
\[
 \frac{\alpha_t}{4t}
 =\frac{2^{t-1}-1}{4t(2^t-1)}
 \le \frac1{24},
 \qquad t\ge2.
\]
Equality holds at $t=2$.  For $t\ge3$, the inequality
$\alpha_t<1/2$ gives
$\alpha_t/(4t)<1/(8t)\le1/24$.  Increasing $G$ relative to $H$ only reduces
the transferred exponent.  The multicollision theorem is stated for fixed
$t$, so the cited result also does not support a growing-multiplicity limit.

Direct composition of the two cited statements therefore yields an exponent
at most $1/24$, rather than $1-o(1)$.  \Cref{thm:main} rules out an
$N^{1-o(1)}$ lower bound at these parameters: the default relation has
quantum query complexity $O(\sqrt N\log N\log\log N)$.  The remark is not
used elsewhere in \cite{JainLiRobereXun2024}, and the black-box separation
theorems of that paper are unaffected.

\paragraph{The lower bound that the reduction does give.}
The case $t=2$ of the composition is valid and yields, to our knowledge,
the best known quantum lower bound.  Let $A_0$ be a graph on $H$ vertices
with no homogeneous set of order $2\log_2H$, which exists by Erd\H{o}s's
bound $R(k,k)>2^{k/2}$ \cite{Erdos1947}.  For $G=H^8/16$, Lemma~III.1 of
\cite{JainLiRobereXun2024} with $t=2$ maps a function $h\colon[G]\to[H]$ to
the graph $A_h$ above on $G$ vertices, and every homogeneous set of order
$\tfrac12\log_2G=4\log_2H-2$ in $A_h$ contains two vertices with the same
$h$-value: an independent set of $A_h$ has distinct $h$-values and maps to
an independent set of $A_0$, so it has fewer than $2\log_2H$ vertices,
while a clique of $A_h$ maps to a clique of $A_0$ and therefore takes fewer
than $2\log_2H$ values on $4\log_2H-2$ vertices.  Finding a collision in a
random function $[G]\to[H]$ requires $\Omega(H^{1/3})$ quantum queries
\cite{Zhandry2015,LiuZhandry2019}, and a query to $A_h$ costs $O(1)$
queries to $h$.  Hence the default Ramsey relation on $N$ vertices has
quantum query complexity $\Omega(N^{1/24})$.  Together with
\cref{cor:succinct}, its quantum query complexity lies between
$\Omega(N^{1/24})$ and $O(\sqrt N\log N\log\log N)$.
